\section*{Chapter \ref{ch:ml:the-machine-learning-team} --- The Machine-Learning Team}

\begin{solution}{exo:ml:the-machine-learning-team:1}
Researchers own the model and its features (the signal, the fit, the validation evidence); \omterm{def:ml:the-machine-learning-team:roles}{machine-learning engineers} own the
service that runs it (online features, compiled model, monitors); the platform team owns what every model shares (\omterm{def:ml:data-and-feature-stores:store}{feature store},
training, registry, serving, the handoff validator).
\end{solution}

\begin{solution}{exo:ml:the-machine-learning-team:2}
Because someone must answer a \omterm{def:ml:monitoring-and-retraining:performance-alarm}{performance alarm}, decide on retraining or retirement and explain the model to validators and
regulators. A team as owner diffuses all three: the page goes to a list, each member assumes another is acting, and the model runs
unwatched.
\end{solution}

\begin{solution}{exo:ml:the-machine-learning-team:3}
$L=\lambda W=0.8\times5.05=4.04$ models, against the chain's 4.00 (theory: $\rho/(1-\rho)=4$ with $\rho=0.8$).
\end{solution}

\begin{solution}{exo:ml:the-machine-learning-team:4}
A window of 5,000 is a valid definition (a known input, a known aggregation, a positive window), and the offline and online
computations both use it, so they agree with each other. Only the comparison of the computed feature values with research's own
values on the reference session shows that the feature is not the one the model was trained on.
\end{solution}

\begin{solution}{exo:ml:the-machine-learning-team:5}
In an M/M/1 queue the mean wait before service is $\rho/(1-\rho)$ service times: 2.8 at 74\% and 9 at 90\%, a little over three
times longer. The chapter's station has two servers and rework, so the numbers differ, but the shape (a wait that explodes near
saturation) is the same, as \cref{fig:ml:team:rate} shows.
\end{solution}

\begin{solution}{exo:ml:the-machine-learning-team:6}
Validation kills 7.9\% of models; review kills 41\%, and the time is lost in queues and rework at the engineers, not in
validation's verdicts. Relaxing validators would let more bad models into paper trading and production without shortening the
lead time much; removing the rewrite (the rework loop between engineers and validators) is what shortens it.
\end{solution}

\begin{solution}{exo:ml:the-machine-learning-team:7}
\texttt{ml\_team.third\_engineer}: with three engineers the engineers' utilisation falls to 50\% and the median lead time to 16.6
weeks (means over three seeds), against 22.3 with two engineers and 11.1 with the package, over the same seeds. Headcount removes
part of the queueing but none of the rework; the package removes both, and frees engineers' time instead of adding to it. The
package first; a third engineer only if the queue returns as research grows.
\end{solution}

\begin{solution}{exo:ml:the-machine-learning-team:8}
Artefact: the coefficient vector and intercept, with their content hash. Features: five definitions with units, sources and the
standardisation fitted on the 60 training days (means and standard deviations as part of the artefact). Feature vectors: raw
inputs and standardised values on reference days. Test vectors: 200 feature rows with predictions to $10^{-12}$. Latency: trivial,
but stated. Monitoring: the five monitors of chapter~27 with their thresholds (0.115, 0.115, 0.089, 0.07 and the CUSUM's $k$ and
$h$), one false page a month, owners. \omterm{def:ml:interpretability-and-model-governance:card}{Model card}: purpose, target (next-day return), horizon, the trading threshold, limitations
(linear, sensitive to units). Checks: as in the chapter, plus a range check on each raw input, which would have caught the
Tuesday of chapter~27 before the model saw it.
\end{solution}

\begin{solution}{pb:ml:the-machine-learning-team:1}
\textbf{Part I.}
\begin{enumerate}
\item Researchers, \omterm{def:ml:the-machine-learning-team:roles}{machine-learning engineers}, platform engineers, the model validator, the trader or portfolio manager, and the
\omterm{def:ml:the-machine-learning-team:roles}{model owner} (a role held by one of them); Kreuzberger and co-authors list seven roles around MLOps.
\item Performance, monitoring, retraining and retirement, and explaining the model to validation, risk and regulators.
\item Everything every model shares (stores, training, tracking, registry, serving, monitors, the handoff validator); not the
models or their features.
\item Because the split determines what passes between teams: a rewrite creates queues and rework, a shared codebase with a
checked package does not; models are harder to modularise than software (Amershi and co-authors).
\end{enumerate}
\textbf{Part II.}
\begin{enumerate}[resume]
\item The forest and its hash, 16 \omterm{def:ml:data-and-feature-stores:store}{feature definitions}, research's feature values on a reference session, 200 test vectors to
$10^{-12}$, a 5-microsecond budget with 2.6 measured, four monitors with owners, a fifteen-field \omterm{def:ml:interpretability-and-model-governance:card}{model card}.
\item Schema, artefact hash, feature specification, feature parity, feature vectors, test vectors, latency, monitoring, \omterm{def:ml:interpretability-and-model-governance:card}{model card}.
\item \Cref{tab:ml:team:defects}: rebuilt artefact, hash and test vectors; unknown input and missing feature, the specification;
millisecond window, feature vectors; another model's outputs, test vectors; library latency, latency; ownerless monitor,
monitoring; card without limitations, the card.
\item A wrong feature can be computed identically offline and online; only a comparison with research's values catches it.
\end{enumerate}
\textbf{Part III.}
\begin{enumerate}[resume]
\item Poisson arrivals at 0.35 a week; review (kills 35\%, 15\% back to research), reimplementation by two engineers (30\% back
for clarification), validation (kills 10\%, 25\% back to the engineers), four weeks of paper trading (kills 20\%), a two-week
canary (kills 5\%).
\item Against the M/M/1 queue: 5.05 weeks simulated against 5, 4.00 models in the system from Book~4's chain, 4.04 by Little's law;
on the whole pipeline 4.87 against 4.90.
\item Most at review (41\%); the time goes to rework loops and to the queue at the engineers, busy 74\% of the time, each model
passing through them 1.79 times.
\item With the rewrite the median climbs from about 16 weeks at 0.2 models a week to 70 at 0.44; with the package it stays near 11.
\end{enumerate}
\textbf{Part IV.}
\begin{enumerate}[resume]
\item \emph{Who owns the model?} With a rewrite at the handoff, the median model reaches production 20.4 weeks after research
(90th percentile 42.9), and 35\% get there: 41\% are lost at review, 5\% at reimplementation, 8\% at validation, 9\% in paper
trading and 2\% in the canary. A validated package instead of the rewrite cuts the median to 11.1 weeks, a saving of 9.3 weeks.
\item Because the gates judge the models, and the models have not changed; only the time spent between gates should.
\item A week of research time per model; it buys 9.3 weeks of lead time and cuts the engineers' load from 74\% to 12\%.
\item No: 16.6 weeks against 11.1, since it shortens the queue but keeps the rework.
\item Median and 90th-percentile lead time, losses per gate, utilisation per team, rework rates per handoff.
\item The researcher who built it or the head of the desk that trades it, named in the registry, not the platform team.
\item It goes back to its author with the failing check's detail, and nothing reaches production until it passes.
\item One codebase for research and production and a package a machine can check.
\end{enumerate}
\end{solution}

\begin{solution}{iq:ml:the-machine-learning-team:1}
Builds and runs the production path: online features equal to offline ones, a model inside its latency budget, deployment,
monitoring, retraining pipelines and rollbacks. The researcher finds and validates the signal; the engineer makes it run without
its author.
\iqlookfor{production concerns: parity, latency, monitoring, operation.}
\end{solution}

\begin{solution}{iq:ml:the-machine-learning-team:2}
The artefact and its hash, the \omterm{def:ml:data-and-feature-stores:store}{feature definitions} in the shared format with reference values, test vectors of inputs and
outputs, a latency budget and measurement, the monitoring specification, a \omterm{def:ml:interpretability-and-model-governance:card}{model card} and an owner, all checkable automatically.
\iqlookfor{a checkable package, not a notebook.}
\end{solution}

\begin{solution}{iq:ml:the-machine-learning-team:3}
Compare offline and online values at the same decision times on a sample (parity) and against research's reference values;
find the first diverging feature and input. Prevent it with one \omterm{def:ml:data-and-feature-stores:store}{feature definition} executed by both paths (a \omterm{def:ml:data-and-feature-stores:store}{feature store}),
parity checks in continuous integration and in production, and feature test vectors in every handoff.
\iqlookfor{parity testing and a single definition.}
\end{solution}

\begin{solution}{iq:ml:the-machine-learning-team:4}
Record every model's timestamps at each stage, and measure waiting against working time, rework loops and the utilisation of each
team; model the pipeline as queues to test changes. Usually the time is in queues before a saturated team and in rework caused by
the handoff.
\iqlookfor{measurement per stage, queues and rework.}
\end{solution}

\begin{solution}{iq:ml:the-machine-learning-team:5}
The platform: storage and feature computation, training infrastructure, tracking and registry, compilation and serving,
monitoring, handoff checks. The model's team: the model, its features and labels, its validation evidence, its monitoring
thresholds and its ownership.
\iqlookfor{shared infrastructure against model-specific choices.}
\end{solution}

\begin{solution}{iq:ml:the-machine-learning-team:6}
The utilisation of each team and the rework rates today: doubling arrivals at a team already busy three-quarters of the time
makes lead times explode long before it runs out of hours. Then decide what to automate (checks, packaging) before hiring.
\iqlookfor{utilisation and queueing near saturation.}
\end{solution}
